\section{Concept: A Perceptual Computer for Safety Evaluation}
\label{sec:concept}

This section derives the system that the study evaluates.
It instantiates the three Per-C components of Section~\ref{sec:bg-cww} for two inputs: the appropriateness scores that
LLM judges give to a reply, and the crisis labels that raters give to the message the reply answers.
Figure~\ref{fig:pipeline} shows the resulting pipeline; the formulas follow the project's method note, and the
implementation is a numpy-only library described in Appendix~\ref{app:prototype}.

\begin{figure}[t]
\centering
\resizebox{\linewidth}{!}{%
\begin{tikzpicture}[
  node distance=5mm and 7mm,
  box/.style={draw, rounded corners=2pt, align=center, font=\footnotesize, minimum height=11mm, text width=27mm},
  io/.style={box, fill=black!6},
  stage/.style={font=\footnotesize\bfseries},
  arr/.style={-{Stealth[length=2mm]}}]
  \node[io] (runs) {Judge runs\\ per reply};
  \node[io, below=of runs] (labels) {Crisis labels\\ per message};
  \node[box, right=of runs] (enc) {IAA panel encoder\\ (IT2 set)};
  \node[box, right=of labels] (deg) {Crisis degree\\ $\mu = k/n$};
  \node[box, right=of enc] (eng) {CWW engine\\ panel / OWA / LWA};
  \node[box, right=of eng] (dec) {Decoder: codebook,\\ Jaccard, 2-tuple,\\ contested flag};
  \node[box, below=of dec] (rules) {Max--min rule base\\ crisis level $\times$ reply word};
  \node[io, right=of rules] (act) {Action\\ A0--A3};
  \node[stage, above=1mm of enc] {Encoder};
  \node[stage, above=1mm of eng] {Engine};
  \node[stage, above=1mm of dec] {Decoder};
  \draw[arr] (runs) -- (enc);
  \draw[arr] (labels) -- (deg);
  \draw[arr] (enc) -- (eng);
  \draw[arr] (eng) -- (dec);
  \draw[arr] (dec) -- (rules);
  \draw[arr] (deg) -- (rules);
  \draw[arr] (rules) -- (act);
\end{tikzpicture}}
\caption{The proposed \Perc{} for safety evaluation. The encoder keeps the disagreement between judges, runs and raters
as fuzzy sets; the engine aggregates the judges; the decoder returns a word, a 2-tuple and a flag for contested
verdicts; a rule base maps the crisis level and the reply word to one of four public-service actions.}
\label{fig:pipeline}
\end{figure}

\subsection{Encoder}
\label{sec:concept-encoder}

All sets are sampled on the universe $X = [0, 6]$ with step $0.01$, which contains the five-point appropriateness scale
with a margin.
A single score $s$ is treated as the Likert bin $[s-0.5, s+0.5]$, the interval of values that round to $s$.

\paragraph{Panel encoder.}
Consider a panel of $J$ judges in which judge $j$ produced runs $r_{j1},\dots,r_{jm_j}$ for the same reply.
Following the IAA \citep{wagner2015iaa}, the encoder builds an IT2 FS whose UMF counts judges and whose LMF counts runs:
\begin{equation}
  \overline{\mu}(x) = \frac{1}{J}\sum_{j=1}^{J} \mathbf{1}\!\left[x \in I_j\right], \qquad
  I_j = \bigl[\min_r r_{jr} - 0.5,\ \max_r r_{jr} + 0.5\bigr],
\end{equation}
\begin{equation}
  \underline{\mu}(x) = \frac{1}{J}\sum_{j=1}^{J} \frac{1}{m_j}\sum_{r=1}^{m_j} \mathbf{1}\!\left[x \in [r_{jr}-0.5,\ r_{jr}+0.5]\right].
\end{equation}
Every run bin of judge $j$ lies inside $I_j$, so $\underline{\mu} \le \overline{\mu}$ holds by construction, and the
number of runs may differ between judges.
The UMF expresses where at least some judges placed the reply; the LMF expresses how consistently each judge did so.
When every rater gives a single score, as the human experts do, the LMF equals the UMF and the set reduces to the
type-1 IAA of the score bins.

\paragraph{Crisis degree.}
For the crisis input, each rater gives a crisp label.
If $k$ of the $n$ raters who gave a label chose a crisis category, the \emph{crisis degree} is $\mu = k/n$, which is the
IAA membership of ``crisis'' \citep{wagner2015iaa}.
Missing labels are not counted in $n$ (decision D13 in Appendix~\ref{app:decisions}).
The crisp alternative is the majority label: crisis if $\mu > 1/2$, no crisis if $\mu < 1/2$, and a tie otherwise.

\subsection{CWW Engine}
\label{sec:concept-engine}

The engine combines the judges of one reply.
The study compares three CWW engines with crisp aggregation (Section~\ref{sec:design-systems}).

\paragraph{Direct panel decoding.}
The simplest engine passes the panel set of Section~\ref{sec:concept-encoder} directly to the decoder.
It treats all judges equally and keeps their disagreement in the FOU.

\paragraph{Ordered weighted average.}
The OWA of values $a_1,\dots,a_n$ with weights $w_j \ge 0$, $\sum_j w_j = 1$, is
$\mathrm{OWA}_w(a) = \sum_{j} w_j b_j$, where $b_j$ is the $j$-th largest value \citep{yager1988owa}.
Its \emph{orness}, $\frac{1}{n-1}\sum_{j}(n-j)\,w_j$, is $1$ for the maximum, $1/2$ for the mean and $0$ for the
minimum.
We obtain weights with a prescribed orness $o$ by mixing the three reference operators linearly: for $o \ge 1/2$,
$w = (2o-1)\,e_{\max} + (2-2o)\,u$; for $o < 1/2$, $w = (1-2o)\,e_{\min} + 2o\,u$, where $u$ is the uniform weight
vector.
An orness below one half gives more weight to the lowest scores, which is the safety-pessimistic attitude used in the
study (orness $0.25$).

\paragraph{Linguistic weighted average with self-consistency weights.}
The LWA computes $\tilde{Y} = \sum_j \tilde{W}_j \tilde{X}_j / \sum_j \tilde{W}_j$ for IT2 inputs $\tilde{X}_j$ and IT2
weights $\tilde{W}_j$ \citep{wu2007lwa}.
Its UMF is the FWA of the input and weight UMFs, and its LMF is the FWA of the LMFs over $\alpha \in [0, h_{\min}]$,
where $h_{\min}$ is the smallest LMF height.
Each FWA is computed on $51$ $\alpha$-cuts: at every level $\alpha$, with cuts $X_j(\alpha)=[a_j,b_j]$ and
$W_j(\alpha)=[c_j,d_j]$,
\begin{equation}
  y_L(\alpha) = \min_{w_j\in[c_j,d_j]} \frac{\sum_j w_j a_j}{\sum_j w_j}, \qquad
  y_R(\alpha) = \max_{w_j\in[c_j,d_j]} \frac{\sum_j w_j b_j}{\sum_j w_j},
\end{equation}
which the Enhanced Karnik--Mendel (EKM) algorithm solves exactly \citep{liu2008fwa,wu2009ekm}.
Here, $\tilde{X}_j$ is the panel set of judge $j$ alone.
The weight $\tilde{W}_j$ is a reliability word chosen from the judge's own consistency on the reply: \word{High} if all
its runs agree, \word{Medium} if they span one level, and \word{Low} if they span two or more.
The weights therefore use no human data and could be computed in deployment.
Table~\ref{tab:wordmodels} lists the weight words.

\subsection{Decoder}
\label{sec:concept-decoder}

\paragraph{Codebook.}
The codebook contains the five appropriateness levels of the BHH protocol, from \word{Harmful/Unacceptable} to
\word{Fully Appropriate} \citep{arnaizrodriguez2026help}.
Word $k \in \{1,\dots,5\}$ is an IT2 trapezoid with core $[k-0.25, k+0.25]$ and shoulders that reach to
$k \pm 0.75$; the parameter $w$ moves the UMF feet outward and the LMF feet inward by $w/2$, so $w$ is the horizontal
width of the FOU (Table~\ref{tab:wordmodels}).
No survey of clinicians' interpretations of these words exists, so $w$ is a proxy; the study uses $w = 0.25$ and
reports sensitivity to it.

\begin{table}[t]
\caption{Word models used by the decoder, the LWA weights and the escalation rule base. Trapezoids are given as
$(a,b,c,d)$. The appropriateness words use the proxy FOU width $w = 0.25$.}
\label{tab:wordmodels}
% Generated by paper/scripts/export_tables.py from results/. Do not edit by hand.
\begingroup\footnotesize
\begin{tabular}{llll}
\toprule
Vocabulary & Word & UMF $(a,b,c,d)$ & LMF $(a,b,c,d)$ \\
\midrule
Reply appropriateness & Harmful/Unacceptable & (0.125, 0.75, 1.25, 1.875) & (0.375, 0.75, 1.25, 1.625) \\
Reply appropriateness & Inappropriate & (1.125, 1.75, 2.25, 2.875) & (1.375, 1.75, 2.25, 2.625) \\
Reply appropriateness & Partially Appropriate & (2.125, 2.75, 3.25, 3.875) & (2.375, 2.75, 3.25, 3.625) \\
Reply appropriateness & Mostly Appropriate & (3.125, 3.75, 4.25, 4.875) & (3.375, 3.75, 4.25, 4.625) \\
Reply appropriateness & Fully Appropriate & (4.125, 4.75, 5.25, 5.875) & (4.375, 4.75, 5.25, 5.625) \\
\midrule
Judge reliability (LWA weight) & Low & (0, 0, 0.25, 0.5) & (0, 0, 0.15, 0.4) \\
Judge reliability (LWA weight) & Medium & (0.25, 0.5, 0.5, 0.75) & (0.35, 0.5, 0.5, 0.65) \\
Judge reliability (LWA weight) & High & (0.5, 0.75, 1, 1) & (0.6, 0.85, 1, 1) \\
\midrule
Crisis level (T1) & None & (0, 0, 0.1, 0.3) & = UMF \\
Crisis level (T1) & Possible & (0.1, 0.3, 0.4, 0.6) & = UMF \\
Crisis level (T1) & Likely & (0.4, 0.6, 0.7, 0.9) & = UMF \\
Crisis level (T1) & Clear & (0.7, 0.9, 1, 1) & = UMF \\
\bottomrule
\end{tabular}
\endgroup

\end{table}

\paragraph{Similarity and word.}
The decoder compares the engine output $\tilde{Y}$ with every codebook word $\tilde{W}_k$ by the IT2 Jaccard
similarity \citep{wu2009comparative}:
\begin{equation}
  s_J(\tilde{Y},\tilde{W}_k) =
  \frac{\sum_x \min\bigl(\overline{\mu}_{\tilde{Y}},\overline{\mu}_{\tilde{W}_k}\bigr) + \min\bigl(\underline{\mu}_{\tilde{Y}},\underline{\mu}_{\tilde{W}_k}\bigr)}
       {\sum_x \max\bigl(\overline{\mu}_{\tilde{Y}},\overline{\mu}_{\tilde{W}_k}\bigr) + \max\bigl(\underline{\mu}_{\tilde{Y}},\underline{\mu}_{\tilde{W}_k}\bigr)},
\end{equation}
and returns the word with the largest similarity.
It also returns the centroid $[c_l, c_r]$ of $\tilde{Y}$, computed with EKM \citep{karnik2001centroid,wu2009ekm}; its
centre $(c_l+c_r)/2$ is a numerical verdict and its width $c_r-c_l$ is an uncertainty measure.

\paragraph{2-tuple.}
For crisp aggregates the decoder uses the 2-tuple representation \citep{herrera2000twotuple}.
A value $\beta$ on the index scale $\{0,\dots,g\}$ of the words maps to
$\Delta(\beta) = (s_i, \alpha)$ with $i = \operatorname{round}(\beta)$ and $\alpha = \beta - i \in [-0.5, 0.5)$, where
halves round up.
The 2-tuple keeps the distance to the nearest word, so that, for example, a mean just below \word{Partially Appropriate}
is reported with a negative offset instead of being silently rounded.

\paragraph{Contested flag.}
A decoded verdict is \emph{contested} if no word fits well or if two words fit almost equally well: the largest
similarity is below $0.5$, or the gap between the two largest similarities is below $0.05$.
A contested verdict can be routed to a stricter policy.
The \emph{escalate-when-unsure} variant flags a reply as unsafe if its word is \word{Harmful/Unacceptable} or
\word{Inappropriate}, or if it is contested and one of these two words is within $0.05$ of the largest similarity.

\subsection{From Verdicts to Public-Service Actions}
\label{sec:concept-rules}

The last component turns two Per-C outputs into an action for a crisis service.
It is a design proposal and has not been clinically validated.

The crisis degree $\mu$ is described by four type-1 words, \word{None}, \word{Possible}, \word{Likely} and
\word{Clear}, whose trapezoids cover $[0, 1]$ (Table~\ref{tab:wordmodels}).
The reply is described by its similarities to the five appropriateness words.
The possible actions, in increasing severity, are:
A0, continue the conversation;
A1, signpost a helpline (in Switzerland 143, or 147 for young people);
A2, request human review or hand over to a counsellor;
A3, urgent human handover and display of the emergency numbers 144 and 112.
Table~\ref{tab:rules} gives the rule table.
It escalates more as the crisis level rises and as the reply becomes less appropriate.

\begin{table}[t]
\caption{Rule table of the escalation decoder (crisis level $\times$ reply word $\to$ action). A design proposal, not
clinically validated.}
\label{tab:rules}
% Generated by paper/scripts/export_tables.py from results/. Do not edit by hand.
\begingroup
\begin{tabular}{lccccc}
\toprule
Crisis level $\backslash$ reply word & Harmful & Inappropriate & Partially & Mostly & Fully \\
\midrule
None & A2 & A2 & A0 & A0 & A0 \\
Possible & A2 & A2 & A1 & A1 & A0 \\
Likely & A3 & A3 & A2 & A1 & A1 \\
Clear & A3 & A3 & A3 & A2 & A1 \\
\bottomrule
\end{tabular}
\endgroup

\end{table}

Inference uses max--min firing.
The strength of the rule for crisis word $c$ and reply word $r$ is $\min\bigl(\mu_c(\mu), s_r\bigr)$, where
$\mu_c(\mu)$ is the membership of the crisis degree in word $c$ and $s_r$ the similarity of the reply to word $r$; the
strength of an action is the maximum over the rules that conclude it.
The decoder outputs the strongest action, and ties go to the more severe action, so that uncertainty never lowers the
response.
With crisp inputs, where a single crisis word and a single reply word have membership one, the inference reduces to a
lookup in Table~\ref{tab:rules}.
